\RequirePackage{etex}

% Vi benytter 2 sidet, a4 størrelse og dansk sprog. 
\documentclass[11pt,a4paper,oneside,danish]{memoir}
\usepackage[paper=A4,pagesize]{typearea}
 \usepackage[utf8]{inputenc} % danske tegn

% Choose language by uncommenting danish or english
\usepackage[danish]{babel} % dansk sprogpakke
% \usepackage[english]{babel}


%% Pakker til fonts, farver, referencer og ligenende er placeret relevant sted i dokumentet.
     
     %% FONT
     \usepackage[default,osfigures,scale=0.95]{opensans} % Alternatively use the option 'defaultsans' instead of 'default' to replace the sans serif font only.
     \usepackage[T1]{fontenc}
     
     % Common packages 
     \usepackage{footnote} % Package used to controle footnote(specially in tables etc.)
     \usepackage{amsmath}
     \usepackage{slantsc}
     \usepackage{multirow} % Merge table rows
     \usepackage{longtable} % Table that can span multiple pages
     \usepackage{geometry}
     \usepackage[section]{placeins} % \FloatBarrier
     \usepackage[compatibility=false]{caption}
     \captionsetup{font=small,labelfont=bf,textfont=it}
     \usepackage{graphicx}
     \usepackage{afterpage}
     \usepackage[draft]{fixme} %hvor 'options' bl.a. inkluderer final, draft, danish, inline og margin
     \usepackage{rotating}
     \usepackage{eso-pic} % til inds�ttelse af forside.
     \usepackage{rotating} % rotate stuff(EVERYTHING!11)

	% % MATH STUFF 
    %\usepackage{algorithmic}
    %\usepackage{harpoon} %Tilf�jer harpun-pile over st�rrelser i math, m.m.
    %\usepackage{pst-plot} % For axes   
    %\usepackage{breqn}
    %\usepackage{fourier}  %................... Roman+math - Utopia
    
  
    %\usepackage{nomentbl} % http://ctan.mackichan.com/macros/latex/contrib/nomentbl/nomentbl.pdf

    %\usepackage{psfrag} % insæt tag, mest brugt i forhold til grafer

%%% DIVERSE CUSTOMS PAKKET NED I FILER FOR OVERBLIK
   % \input{styles/HaxTilPreamble.tex}
   % \input{styles/CustomFunctions.tex}
%	\input{Styles/code.tex} 
    \usepackage{todonotes}


% new commands for individually colored todo 
\newcommand{\todosander}[2]{
\todo[#1, color=green!40]{#2 \\- Sander}
}
\newcommand{\todosimon}[2]{
\todo[#1, color=red!40]{#2 \\- Simon}
}
\newcommand{\todorasmus}[2]{
\todo[#1, color=yellow!40]{#2 \\- Rasmus}
}
\newcommand{\todokristoffer}[2]{
\todo[#1, color=purple!40]{#2 \\- Simon}
}


% Color package. Use with listings package and use for color a table
%\usepackage[usenames,dvipsnames,svgnames,table]{xcolor}

    %%%%%%%%%%%%%%%%%%%%%%%%%%%%%%%%%%%%%%%%%%%%%%%%%%%%%%%%%%%%%%%%%%%%%%%%%%%%%
%                            Custom for funktioner                          %
%                     			      	        		                    %
%%%%%%%%%%%%%%%%%%%%%%%%%%%%%%%%%%%%%%%%%%%%%%%%%%%%%%%%%%%%%%%%%%%%%%%%%%%%%


% Package for wrapping figures in text
\usepackage{wrapfig}

% Set the path where LaTeX looks for pictures.
\graphicspath{{figure/}}

% You need a newsubfloat element to use subcaption
\newsubfloat{figure}

% Command to set caption styles
\captionnamefont{\bfseries\small}
\captiontitlefont{\itshape\small}
\subcaptionlabelfont{\bfseries\small}
\subcaptionfont{\itshape\small}

% When we don't use chapter numbering, 
% the figure numbering are shown like 
% fx 0.3 if we don't use this line.
\counterwithout{figure}{chapter}


% Command \figur{filename}{caption}{label}{width} 
% for inserting a new figure.
\newcommand{\figur}[4]{ 
 \begin{figure}[ht] 
 \centering 
 \includegraphics[width=#4\textwidth]{#1} 
 \caption{#2} 
 \label{#3} 
 \end{figure}
 } 
 
% Command \wrapfigur{filename}{caption}{label}{width}{r/l}
% for inserting a new figure wraped in the text. 
\newcommand{\wrapfigur}[5]{ 
 \begin{wrapfigure}{#5}{#4 \textwidth}
   \begin{center}
    \includegraphics[width=#4\textwidth]{#1}
   \end{center}
   \caption{#2}
   \label{#3} 
 \end{wrapfigure}
}
 
% Command \subfigur{filename1}{filename2}{caption}
% {subcaption1}{subcaption2}{label1}{label2}
% {width1}{width1} for inserting a new figure.
\newcommand{\subfigur}[9]{
\begin{figure}[ht]
\centering
\subbottom[{#4}\label{#6}]%
    {\includegraphics[width=#8\textwidth]{#1}}\hfill
\subbottom[{#5}\label{#7}]%
    {\includegraphics[width=#9\textwidth]{#2}}
\caption{#3}
\label{fig:subfigures}
\end{figure}
}
% Command \truefigure{filename}{caption}{label}
% for inserting a picture without scaling.
\newcommand{\truefigure}[3]{ 
 \begin{figure}[ht] 
 \centering 
 \includegraphics{#1} 
 \caption{#2} 
 \label{#3} 
 \end{figure}
 } 

% Keeps floats in the section. 
% Help control place to put pix
\usepackage[section]{placeins}

% Commands for rightaligning wraped text in a tabular
\newcommand{\rr}{\raggedright}
\newcommand{\tn}{\tabularnewline}

\usepackage{multirow}

%Tabular pakke som Wex har fundet, men som vi ikke bruger alligvel
%\usepackage{tabu}

% Include code pieces from files
% Documentation:
% ftp://ftp.tex.ac.uk/tex-archive/macros/latex/contrib/listings/listings.pdf
\usepackage{listings}


%Package with enumirate command with extra options (used within cells of tabular)
\usepackage[inline]{enumitem}
    
    % add a box using the themecolor. Can be used to wrap around a chapintro for introducing an assignment
    %use \begin{assignment} and end\{assignment} to box it
    %http://texdoc.net/texmf-dist/doc/latex/mdframed/mdframed.pdf
     
    \usepackage[xcolor]{mdframed} 
    \newmdenv[linecolor=ThemeColor]{assignment}
%    \usepackage{Placeins}
    
% % % % % % % % % % % % % % ********% % %Generelt design % % % % % % % % % % % % % % % % % % %
 
%\usepackage[table,dvipsnames]{xcolor} % 'table' er i stedet for \usepackage{colortbl} og 'dvipsnames' er et farveskema hvor man angiver farver med navne

\usepackage{xcolor} % farver kan indsætttes i henhold til RGB med definecolor

%Valg af temafarve % Denne farve bruges i chapterstyle %
\definecolor{ThemeColor}{RGB}{209, 67, 26}

% Valg af font %
\usepackage[sc]{mathpazo}
\linespread{1.05}         % Palatino needs more leading (space between lines)
\usepackage[T1]{fontenc}

% Pakker inkludere overordnede designvalg. Pakkerne er valgt for overskuelighed %
	 %%%%%%%%%%%%%%%%%%%%%%%%%%%%%%%%%%%%%%%%%%%%%%%%%%%%%%%%%%%%%%%%%%%%%%%%%%%%%
%                             Chapter and Section                           %
%                     			      Style	       		                    %
%%%%%%%%%%%%%%%%%%%%%%%%%%%%%%%%%%%%%%%%%%%%%%%%%%%%%%%%%%%%%%%%%%%%%%%%%%%%%


% % % % Redefined so that header will show even on chapter start. % % % %
\makeatletter
    \let\stdchapter\chapter
    \renewcommand*\chapter{%
    \@ifstar{\starchapter}{\@dblarg\nostarchapter}}
    \newcommand*\starchapter[1]{%
        \stdchapter*{#1}
        \thispagestyle{fancy}
        \markboth{\MakeUppercase{#1}}{}
    }
    \def\nostarchapter[#1]#2{%
        \stdchapter[{#1}]{#2}
        \thispagestyle{fancy}
    }
\makeatother




\usepackage{titlesec} % Pakke til opsætning af chapter og sektioner


% % defining new font for the chapternumber. Calling it FuckingHuge 
\newcommand{\FuckingHuge}{% Needs a 'stretchable' font
      \usefont{\encodingdefault}{\rmdefault}{b}{n}%
      \fontsize{35}{45}\normalfont\color{ThemeColor}} %60 80


%%% TAL I OVERSKRIFTER, både chapters og sections. 
    \maxsecnumdepth{subsubsection} %subsubsection er dybdeniveau for numering.
   	\setsecnumdepth{subsubsection}


  
 % % % % CHAPTER OPSÆTNING % % % % % % % % % % % % % % % % % % % % %

%Hang thoose numbers! I margin med tallene for både sections og chapters! En lækker chapter style!
\chapterstyle{hangnum}
 
%Modificering af ovenstående style
%Placering af chapter overskriften på siden i forhold til sidens start %
\renewcommand*{\chapterheadstart}{\vspace*{-20pt}}


% FARVE PÅ CHAPTER % %
%TItel farve og font valg
\renewcommand*{\chaptitlefont}{\normalfont\Huge\color{ThemeColor}}
%Numering farve og font valg
\renewcommand{\chapnumfont}{\normalfont\FuckingHuge\color{ThemeColor}}



% % % % SECTION OPSÆTNING % % % % % % % % % % % % % % % % % % % %
\hangsecnum % vi smidder numering ud i marken

%Vi benytter tema farven til numeringen
\renewcommand\thesection{\color{ThemeColor}\thechapter.\arabic{section}}



    
    


	% \input{Styles/header.tex}
	 
% Setup the margins:
\setlrmarginsandblock{1in}{*}{1.5}
\setulmarginsandblock{3cm}{4cm}{*}

% Setup header and footer: %
% \setheadfoot{headheight}{footskip}
% \setheaderspaces{headdrop}{headsep}{ratio}
\setheadfoot{45pt}{1cm}
\setheaderspaces{1cm}{*}{*}
\checkandfixthelayout
% Se dette pdf for dokmentaton:
% www.tug.dk/_media/foredrag/2003-12-03/slides-src.pdf

 \checkandfixthelayout[nearest]
    \checkthelayout
    \fixthelayout
    \baselineskip \topskip

%BibTex setup
\usepackage[semicolon,authoryear,square]{natbib}

\bibliographystyle{plainnat}
%%%%%%%%%%%%%%%%%%%%%%%%%%%%%%%%%%%%%%%%%%%%%%%%%%%%%%%%%%%%%%%%%%%%%%%%%%%%%%%%%%%%%%%%%%%%%%

% % Smart listing function. Documentation: http://texblog.org/tag/makenomenclature/
\usepackage{nomencl}
\makenomenclature 

% % % % % % % % % % % % % % % % Reference opsætning % % % % % % % % % % % % % % % % % % % % % % %

\usepackage{varioref} % for references, http://www.ctex.org/documents/packages/bibref/varioref.pdf
   
% External reference package, zref makes the references into hyperlinks.
\usepackage{url}
\usepackage{hyperref}
%\usepackage{zref-xr}
\usepackage{xr}

% % % % % % % % % % % % % % % % % % % % % % % % % % % % % % % % % % % % % % % % % % % % % % % % % % %

%\usepackage{showlabels}

% footruleskip is allready defined in memoir, so we have to undefine it, in order to include fancyhdr for our footer and header.
\let\footruleskip\undefined
\usepackage{fancyhdr}% For making fancy headers http://ctan.org/pkg/fancyhdr

\AtBeginDocument{\renewcommand\contentsname{Indholdsfortegnelse}}

% % % % % % % % % % % % % Inkluderinger af ønskede kapitler % % % % % % % % % % % % %

\begin{document}
\newgeometry{left=3cm,bottom=0.1cm}
% Write the title
{
\centerline{\Huge\huge\bfseries\color{ThemeColor}Team 15107} 
\vspace{1em}
\centerline{\large\color{ThemeColor}Aarhus School of Engineering - Aarhus Universitet}
}

\centerline{\large\date{10. April 2015}}

%\color{Black}
%\vspace{8em} % Tilføj nogle tomme linjer
\begin{center}
    \begin{tabular}{ l l p{6cm} }
	 & & \\
	& & \\
    Stud. nr.: 201270115 & Johnny Davidsen Poulsen & \\\hline
    & & \\
    Stud. nr.:  02255 & Kim Rose Faber & \\\hline
    & & \\
    Stud. nr.:  & Kristian Rossen & \\\hline
     \end{tabular}
      	\vspace{50pt}
      	
      \textbf{Repository}\\
     http://gitswat.ase.au.dk/git/15/1/team15107/ex13.git
\end{center}
\thispagestyle{empty} % No header and footer on this page
\restoregeometry
\setcounter{tocdepth}{2}
\tableofcontents

 

% Setup header and footer
\pagestyle{fancy}
\fancyhf{} % Clear all header an footer fields 
\renewcommand{\headrulewidth}{0pt} % Delete headruler
\fancyhead[R]{\includegraphics[height=40pt]{header}} % Insert picture in left side of header
%\fancyhead[R]{\rightmark} % insert chapter title in right side of header
\fancyfoot[LE,RO]{\thepage} % insert pagenumber on right side of even page, left side of odd page 
%\setcounter{page}{1}

\thispagestyle{fancy}
\chapter{Første itteration}
\label{chp:FirstItteration}

\section{Design og implementering}
% Johnny

\figur{sd_tick_firstitteration}{Sekvensdiagram for ATM.Tick kald}{fig:sd_tick_first}{1}

Første itteration af ATM systemet havde en klasse kaldet ATM. ATM klassen havde to metoder; Tick og AddTrack. AddTrack tog en Liste af Tracks som parameter og tilføjede disse Tracks til listen af Tracks som ATM klassen skulle håndtere. 

Et objekt af ATM blev oprettet sammen med en Timer i programmets main metode. Timeren blev sat til et interval på 250ms med en delegate som kaldte Tick hver gang Timeren udløb. Implementeringen af Tick metoden kan ses som et sekvensdiagram i figur \ref{fig:sd_tick_first}. Figuren viser hvordan at Tick Kalder UpdatePositions i TrackHandler med listen af vores tracks. UpdatePositions opdatere Tracksenes position baseret på hastighed, position og tidspunkt. Herefter sender Tick den opdaterede liste til RegionHandleren ved at kalde metodens UpdateRegion metode. UpdateRegion identificerer, håndterer og logger events ved hjælp af EventLoggeren. EventLoggeren har en metode for hver type af event som skal logges.

SOLID principperne er tænkt med i designet således at en ny implementering af ITrackHandleren, IRegionHandleren og IEventLogger kan injectes hvis implementeringen skal ændres. ATM'en burde således være åben for ny implementering af håndtering af tracks og region uden at skulle modificeres.

\section{Test}
Der blev kun skrevet ganske få test i den første itteration inden at kravene til systemet blev opdateret. Nogle få testmetoder blev skrevet som testede for korrekt håndtering af null parametre ved nogle metodekald.


%\part{\section{Overvejlser}
%I udarbejdelsen af designet har SOLID principperne været i fokus. Herunder har specielt Single Responsibility haft stor indfyldelse for at skabe software, som var nemt at unit teste. Dette er gjort med et ønske om at skabe en skarp ansvarsfordeling således at unit testene ville være selvforklarende og præcise.
%
%Herudover har Open/Closed prinippet været vigtige i overvejelserne, da det allerede fra start var annonceret, at der ville komme ændringer til kravene. Der er således forsøgt at lave et design, hvor der nemt kunne tilføjes nye funktioner uden at ændre i allerede skrevet kode. Der er derfor kodet til interfaces istedet for til specifikke klasser for at gøre klasseudskiftning og samtidig også testning lettere.
%
%
%
%\section{Implementering}
%
%I den endelige implementering, som ses i figur \ref{fig:klassediagagram}, blev der valgt en klar opdeling mellem udskrivning af output og  håndtering af tracks og relevante events. Dette blev opbygget med en OutputControl og en RegionControl. OutputControl håndterer at skrive både til konsol og fil. 
%RegionControl har en liste af ITracks og IEventDetectors således, at listen af ITracks ved hvert ticks sendes  til de forskellige EventDetectors, som behandler tracks i forhold til deres event. De respektive Events ved ligeledes selv, hvordan de skal udskrives.
%
%\figur{AirTrafficControlSystem_v3}{UML Klassediagram}{fig:klassediagagram}{1} 
%
%
%\section{Håndtering af kravændringer}
%Overordnet gjorde en grundig brug af OCP det simpelt at tilføje de nye krav til projektet. 
%
%En klar opdeling af events og detektorer gjorde det simpelt at tilføje en AirportDetector samt AiportEvent, som kunne håndtere den nye type af events.
%
%Ligeledes blev konfigurationsfilens parameter blot sendt med som constructorparameter til de respektive klasser, og dette kunne ligeledes gøres med minimale ændringer. NearMissRadious var valgt hard coded i NearMissDetector, og ændringer af dette krævede kun mindre ændringer i allerede skrevet kode.
%
%Overordnet vurderes der derfor en god balance mellem arbejdsmængden og hvor generisk softwaren var, da de nye krav skulle tilføjes.
%}
\thispagestyle{fancy}
\chapter{Anden itteration}
\label{chp:seconditteration}
Efter at første itteration var færdig blev et nyt sæt af krav til systemet præsenteret. På den baggrund blev systemets design opdateret således at det bedre ville kunne udviddes i fremtiden med flere ændringer af samme slags. 

\section{Design og implementering}
De nye krav dikterede at systemet skulle opsættes ved hjælp af en 

\figur{sd_timer_elapsed}{Sekvensdiagram for ATM.\_ticktimer\_Elapsed kald}{fig:sd_timer_elapsed}{1}

*RegionElements
*Configfile
*IEvents
*EventDetectors

\todo[inline]{Johnny}

\section{Unittest}
\todo[inline]{??}
\subsection{ATM}
\todo[inline]{KRK}
\subsection{TrackHandler}
\todo[inline]{Johnny}
\subsection{RegionEventDetector}
\todo[inline]{Johnny}
\subsection{EventDetectors}
\todo[inline]{Kim}
\subsection{Events}
\todo[inline]{Kim}
\subsection{Excluded code}
\todo[inline]{Kim}
\section{Integrationstest}
\todo[inline]{Johnny}

%
%\section{Teststrategi}
%Gruppen valgte at enhedsteste hver metode så snart den var blevet skrevet. Dette sikrede for det første, at der let kunne holde en coverage nær 100\%, da der blot blev sørget for at få fuld coverage hver eneste gang en ny metode var oprettet. Dette gav en god sikkerhedsfølelse da man da altid kunne stole på, at enhederne virkede i en individuel sammenhæng, og at der derfor ikke var langt til en fungerende integration.  Det gav samtidig en rigtig god afveksling mellem udvikling og test, så arbejdet ikke blev for monotomt. 
%Da enhederne i bunden af hierakiet var blevet testet, påbegyndte vi integrationstest nedefra og op, og til sidst blev hele systemet testet i en afsluttende "big bang" integration.
%
%\subsection{Enhedstest}
%Alle enhedstests blev så vidt muligt konstrueret ved hjælp af isolation frameworket NSubstitute, og testet i Unit-test frameworket NUnit. 
%NSubstitute hjalp os til hurtigt at skrive gode, isolerede unittests. Ulempen ved at bruge et isolation framework var, at man hurtigt glemte om man kodede til en stub eller en mock, fordi det med frameworket var det samme. Vi kom derfor let til at bruge mange mocktests, selvom nogle testcases måske godt kunne være klaret med en stubtest. 
%Under udviklingen af unittests kom gruppen også let til at fokusere udelukkende på coverage. Vi oplevede flere gange at måtte gå tilbage og revurdere tests, fordi de reelt set ikke testede det vi ønskede, på trods af at de havde fuld coverage. Dette var vigtig læring, da man ellers hurtigt kunne stirre sig blind på coverage. 
%
%\subsection{Integrationstest}
%Den primære strategi for integrationstest var som nævnt tidligere bottom-op. Vi startede med at teste de nedre klasser i hierakiet, for derefter at bevæge os opad. De nederste klasser gav en del udfordringer, da det kan være rigtig svært at teste systemets afhængigheder til styresystemet osv. Dette gjorde det blandt andet svært for os at teste output til konsollen, som praktisk talt er umuligt med mindre man intercepter en stream. Problemet ved intercepte en stream er, at man da vil skulle teste i uendelighed, da man da også skal teste sin interception osv. Output til filer blev testet ved at skrive til en fil, læse filen ind igen, for så at sammenligne om det rent faktisk var blevet tilføjet til filen. Efterhånden som vi bevægede os højere op i hierakiet, blev det lettere for os at lave integrationstests, da vi da bedre selv kunne styre afhængighederne.
%Da der var blevet integrationstestet hele vejen op istedet for blot at gøre brug af en big bang test, havde vi en klar formodning om, at systemet ville virke som forventet, hvilket det også gjorde.
%
%\section{Test som support ved ændrede krav}
%De ændrede krav resulterede i, at en række tests skulle skrives om. Nogle få af vores ændrede cases kunne være undgået med et bedre design, men langt de fleste ville være nødvendige uanset designet. 
%De ændrede krav betød også at en række constructors i vores klasser skulle skrives om. Her oplevede vi virkelig hvordan en setup metode kan være en fordel, da vi her blot skulle ændre ét sted, hvorefter alle vores tests var opdateret til at supportere de nye ændringer.
%De nye ændringer krævede naturligvis også en række nye testcases, hvilket hurtigt blev tilføjet til projektet. 
\thispagestyle{fancy}
\chapter{Continous Integration}
\label{chp:contintegration}

\section{Versionstyring}
Versionsstyringsværktøjet Git og tilhørende Git server gjorde det muligt, for hvert gruppemedlem at arbejde individuelt med en del af opgaven. Ved afsluttet delopgave blev dette tilføjet til det samlede projekt på Git serveren. De andre gruppemedlemmer skulle pulle og merge deres arbejde med den nyeste udgave a projektet inden deres eget arbejde kunne pushes til serveren. Der kunne til tider opstå merge conficts hvilket resulterede i lidt ekstra arbejde, men da versionsstyring og mergetools er tilgængelig blev disse problemer let løst. Selv om det er muligt, at arbejde på separate brances blev dette ikke særlig aktuelt, da de enkelte arbejdsopgaver ikke var bredtfavnende.   

\section{Buildserver}
Jenkins er benyttet som buildserver til at udelukke alle lokale indspillende faktorer ved compiling af projektet. Jenkins er linket op mod gruppens Git repository og bygger den seneste udgave af projektet hver gang der pushes til branchen master. Efter projektet er blevet bygget af Jenkins rapporterers der tilbage om projektet kunne compiles eller ej. Udfaldet bliver lagt i en byggehistorik, hvor det også er muligt at se detaljer om tidligere byg. Alle medlemmer af gruppen fik garanteret, at deres tilføjelser til det samlede projekt ikke spolerede udviklingsarbejdet hos de andre i gruppen når der skulle pulles og merges. Endvidere gjorde to plugins til DotCover og NUnit det muligt, at se om alle tests og codecoverage var som forventet.

\subsection{DotCover}
Dot cover er en statisk kodeanalyse værktøj der analyserer hvor mange linjer kode vi tester.  Det er muligt at ekskludere code der ikke skal tages med i analysen. Det kan være klasser, moduler eller funktioner. Vi har besluttet at code coveres skal være 100 \%, det betyder dog ikke at vi kan være sikre på at kodens funktionalitet er testet, men blot at alle linjer kode har været berørt af vores unit test. Det vil sige at det ikke er funktionaliteten der er testet.

\subsection{NUnit}
NUnit er et open source testing framework for udvikling af projekter i .Net. Alle unittests er baseret på NUnit konventioner og gruppen benyttede sig af dette til at teste klasser og underliggende metoder. Disse tests blev også eksekveret på Jenkis serveren, og resultatet der fra vist på projektets forside.

\figur{NUnit_test_results}{NUnit testresultater fra Jenkins}{fig:NUnit_test_results}{0.8}

%\section{Strategi}
%Fra start var gruppen enig om at dele sig ud i 2 par og arbejde i pair programming. På den måde kunne der arbejdes på to dele af designet på en gang. Kvaliteten blev dermed højnet af at der var fire øjne på skærmen, mens der blev programmeret og taget beslutninger herom.\newline
%ControlSystem blev delt op i to forgreninger, nemlig RegionControl og OutputControl. Hvert par startede fra bunden på hver sin forgrening. Målet med at arbejde på denne måde er at de to par kan implementere og teste hver deres forgrening, hvorefter det til sidst kan stykkes sammen. Undervejs blev der lavet lidt om i designet og der kom flere opgaver til i RegionControl. En ny opgave blev blot taget af et par, når en opgave var færdig. 
%
%\section{Continous Integration}
%Under udviklingen af ATC blev der gjort brug af Continous Integration vha. git repository og Jenkins. Et startprojekt er blevet oprettet, og i det, har hvert par udviklet hver deres kode, klasser og tests. Løbende blev der committet og pushed til git serveren, så gruppen nemt kunne følge med i hvor langt det andet par var i udviklingen. På denne måde blev der dannet en versionshistorik og det var muligt, hvis alt gik galt, at gå tilbage til en tidligere version, hvor koden virker.\newline
%
%
%\subsection*{Jenkins}
%Jenkins bruges til at følge med i hvert byg, der bliver committet til git. Her bliver det automatisk testet med de unit- og integrationstests, der er skrevet til hver klasse. Jenkins viser en byggehistorik, og ud fra de fem seneste byg laves der en beskrivelse i form af en sol hvis de er gået godt. Derimod er der en grå sky med lyn, hvis de seneste fem byg er fejlet. En coverage report kan også ses i Jenkins, og alle kan følge med i hvor meget af koden, der er testet. Jenkins genererer et konsolouput, så der kan aflæses, hvor der kan være eventuelle fejl i koden.\newline
%Når man i gruppen kender hinanden godt, er Jenkins et godt og motiverende redskab. Som udvikler og tester er man ikke interesseret i at ens byg og test fejler, og derfor vil man gerne lave et ordentligt stykke arbejde, så det kan pushes til git og Jenkins uden fejl. 
%
%\begin{figure}[htpb]
%    \centering
%    \includegraphics[width=0.7\textwidth]{JenkinsGraf}
%    \caption{Graf fra Jenkins ober tests og byg}
%    \label{fig:JenkinsGraf}
%\end{figure}
%\FloatBarrier
%%\figur{JenkinsGraf}{Graf fra Jenkins over tests og byg}{fig:JenkinsGraf}{0.7} 
%
%På grafen ses ud af x-aksen de antal byg Jenkins har lavet. Op ad y-aksen ses antal test cases der er skrevet. At grafen er blå, viser at der ikke er nogen tests der har fejlet når jenkins har bygget.
%
%\section{Udfordringer}
%Under forløbet med udvikling af ATC opstod forskellige udfordringer. Herunder var det en udfordring af få git til køre som ønsket. Hertil skal det være stabilt, og alle i gruppen skal pushe, gerne uden konflikter. På den anden side skal de andre pull hvor alt lykkedes. Når der committes/pushes/pulles er der adskillige problemer, der kan opstå når de forskellige filer skal merges. Blandt oplevede gruppen en udfordring med at merge, hvilket af flere gange gav følgende fejl: \textquotedblleft You can't merge, because you have unmerged files \textquotedblright. I denne situation blev der ikke fundet anden udvej end at gemme sine file andetsted lokalt, og clone projektet igen. Heri indsættes de gemte filer på my, og merge af projektet lykkedes.

\thispagestyle{fancy}
\chapter{Konklusion}
\label{chp:Konklusion}
Der har gennem projektperioden vist sig et behov for at teste. Flere tests har afsløret væsentlige fejl hvor der eksempelvist blev ittereret gennem lister i forkert rækkefølge eller beregninger blev returneret med værdier som havde forkerte værdier på 12. decimal. Anden itteration gav et design som er meget nemmere at udvidde med lignende krav. Fordelen ved at implementere et design som bygger på SOLID principperne har vist sine fordele. Vi har gennem dette projekt udviklet en bedre forståelse for SOLID og relevansen af unittesting. Integrationenstestene er desværre ikke implementeret, men der er tilgengæld lavet en række test ved hjælp af et eksekverbar projekt hvor en række scenarier er testet manuelt. Dette har givet os en god fornemmelse af at projektet fungere uden at vi dog har automatiseret processen.

%I forbindelse med opgaven har gruppen fået en bedre forståelse for vigtigheden af en godt og testbart design med en klar opdeling af klasserne. Dette gjorde det også forholdsvist smertefrit at håndtere de ændrede krav, der kom sent i forløbet, og det har givet gruppen et indblik i, hvad man kan blive udsat for i et virkeligt projekt. 
%
%Brugen af Jenkins var en motiverende faktor for at lave tests løbende, da man kan følge med i grafens udvikling, og det har haft stor betydning for udviklingsforløbet at projektet. 
%
%Ydermere har gruppen fået et væsentligt bedre erfaring til brugen af testframeworks i form af NUnit, DotCover og NSubstitute. Både i forhold til det rent tekniske, men i samme grad at skabe en god struktur og præcision i de udarbejdede tests.


%\bibliographystyle{plainnat}
%\bibliography{bibliografi}% Selects .bib file AND prints bibliography
\thispagestyle{fancy}

%\listoftodos % Make a list of todo's
\end{document}